A tiny model is given one job. It reads a sentence --- \emph{``write a
function that takes two numbers, a and b, and returns their sum''} ---
and it writes back the right Python:

\begin{Shaded}
\begin{Highlighting}[]
\KeywordTok{def}\NormalTok{ solve(a, b):}
    \ControlFlowTok{return}\NormalTok{ a }\OperatorTok{+}\NormalTok{ b}
\end{Highlighting}
\end{Shaded}

We run the hidden tests. It passes. If this were a demo, we'd move on.

But this is not a demo. So before we celebrate, we ask the one question
that turns this whole book into an experiment rather than a party trick:

\begin{quote}
\textbf{Why did it get it right?}
\end{quote}

Maybe it genuinely learned that ``sum'' and ``a and b'' mean ``add them
together.'' Maybe it memorized that exact sentence from training. Maybe
it noticed this was nearly identical to an example it saw and copied the
answer. Maybe a much dumber machine --- one that does nothing but look
up the closest training example --- could have done the same thing. Or
maybe our test wasn't as hard as we thought.

Making the model answer is easy. \textbf{Knowing why it answered is the
whole game.}

\section{One job, and a way to check
it}\label{one-job-and-a-way-to-check-it}

We give the model exactly one, narrow, checkable job:

\begin{quote}
\textbf{read a plain-English programming request → produce a working
Python program}
\end{quote}

Then we \emph{run} that program against hidden tests. That single choice
--- executable output plus hidden tests --- is what makes this a real
experiment instead of a showcase, because it hands us the things a
language-model toy usually lacks:

\begin{itemize}
\tightlist
\item
  inputs we control,
\item
  code we can actually execute,
\item
  an objective pass/fail,
\item
  data we can generate in any quantity,
\item
  difficulty we can dial up or down,
\item
  tests we can keep hidden from the model,
\item
  splits we can design to expose memorization,
\item
  audits we can run to catch shortcuts.
\end{itemize}

\section{What this book is (and
isn't)}\label{what-this-book-is-and-isnt}

This is not a transformer tutorial. It is not a guide to chatbots,
agents, or coding assistants. It is a careful, honest investigation of a
single question:

\begin{quote}
\textbf{How much structured programming capability can a small language
model actually learn --- and how do we tell whether it \emph{learned}
it, rather than memorized it or found a cheap shortcut?}
\end{quote}

\section{The part that's actually
hard}\label{the-part-thats-actually-hard}

The natural guess is that \emph{building} a small model is the hard
part. It isn't. Once the laboratory works, a \textasciitilde0.9
million-parameter model will reliably turn a handful of programming
sentences into correct Python --- in-distribution, that is nearly easy.

The hard part is everything after: proving it didn't just copy, checking
whether it still works when the \emph{words} change but not the
\emph{computation}, testing whether it can \emph{combine} two things it
knows into a third thing it's never seen, and auditing our own evaluator
to make sure the number it reports means anything.

This book builds a laboratory, runs those tests, and reports --- plainly
--- what happened. Including the moments the model fell flat, and the
moment we discovered our own measuring stick was broken, not the model.

\textbf{The model was the instrument. The experiment was the subject.
Let's brew.}
